% !TEX root = ../notes_missingsemester.tex
% Block 10: AI Dev & Vibe Coding
%%%%%%%%%%%%%%%%%%%%%%%%%%%%%%%%%%%%%%%%%%%%%%%%%%%%%%%%%%%%%%%%%%%%%%

\index{AI-assisted development}\index{vibe coding}\index{LLM}\index{agent}
\chapter{AI Dev \& Vibe Coding}
\printchapterversion{10_ai_dev_vibecoding}
\newthought{The tools changed. The fundamentals did not. Know both.}

% ==============================================================================================
% BLOCK 1: INTRODUCTION
% ==============================================================================================

\section{Two Ways AI Helps You Code}

\newthought{AI coding tools} come in two flavours, and knowing when to use which matters:

\begin{description}
    \item[Terminal/CLI agents]\index{CLI agent} Claude Code\index{Claude Code} (\url{https://docs.anthropic.com/en/docs/claude-code}),
    Aider\index{Aider} (\url{https://aider.chat/}),
    Codex CLI work directly in your repo from the terminal, read your files, make changes, run commands.
    Best for: implementing features, debugging, refactoring across multiple files.
    \item[IDE-integrated agents]\index{IDE agent} Cursor\index{Cursor} (\url{https://www.cursor.com/}),
    Windsurf (\url{https://windsurf.com/}), Cline (\url{https://cline.bot/}) full IDE with AI
    built in, visual diffs, inline suggestions. Best for: editing in context, reviewing changes visually,
    inline completions while typing.
\end{description}

\marginnote{Autocomplete tools: GitHub Copilot\index{GitHub Copilot} (\url{https://github.com/features/copilot}).
Chat-based tools: ChatGPT (\url{https://chatgpt.com/}), Claude (\url{https://claude.ai/}).
CLI and IDE agents are the next level: they read your whole project, not just the current file.}

% ==============================================================================================
\newpage
\section{Optional Exercise: Catch Up to \texttt{v0.8}}\index{exercises}\index{Git!tag}\index{recovery}

\marginnote{\newthought{Tag map.}\index{Git!tag}
\texttt{v0.8} marks the end of Block~9.
This is the final block, no further tag is produced.
Browse the full set at
\url{https://github.com/schutera/pixelwise/tags}.}

\newthought{Catching up via tag.}\index{Git!tag}\index{recovery}
This catch-up extends the one in Block~9. If you have not done that one,
do it first: It lands you at \texttt{v0.7}, the end-of-Block~8 state.
The exercise here picks up from there and rewinds to \texttt{v0.8}, the
end-of-Block~9 state, with structured JSON logging, an upgraded
\texttt{/health} endpoint, a Discord webhook for error-level alerts,
and the \texttt{/api/analytics} endpoint serving aggregates from the
PostgreSQL \texttt{predictions} table.

\newthought{Both VMs this time.} Block~9 added observability and
analytics on top of the running service, so this catch-up brings both
machines up to \texttt{v0.8}. The Git-tracked changes since \texttt{v0.7}
include structured JSON logging in \texttt{app/main.py} that emits
timestamps, levels, and per-request fields, an upgraded \texttt{/health}
endpoint that distinguishes liveness from readiness,
\texttt{scripts/health\_check.sh} for the cron-driven probe, the Discord
webhook integration that posts on error-level logs, the
\texttt{/api/analytics} endpoint reading aggregates from the
\texttt{predictions} table, and a refreshed \texttt{.env.example} with
\texttt{DISCORD\_WEBHOOK\_URL}. The local \texttt{.env} on each VM needs
the new webhook URL written by hand; \texttt{setup-server.sh} then
reconciles dependencies and restarts the service.

\newthought{Run the catch-up on \texttt{dev}.}
Rewind the repository to \texttt{v0.8}, refresh \texttt{.env} from the
updated example and edit it to set \texttt{DISCORD\_WEBHOOK\_URL},
reinstall pinned dependencies if your virtual environment was wiped,
then run \texttt{setup-server.sh}. The exact commands are in the margin.

\marginnote{\newthought{Sequence on \texttt{dev}.}\\[0.4em]
\texttt{cd \textasciitilde/pixelwise}\\
\texttt{git fetch origin -{}-tags}\\
\texttt{git reset -{}-hard v0.8}\\
\texttt{cp .env.example .env}\\
\texttt{nano .env}\\
\texttt{source .venv/bin/activate}\\
\texttt{pip install -r requirements.txt}\\
\texttt{bash setup-server.sh}\\[0.3em]
Edit \texttt{.env} to set \texttt{DISCORD\_WEBHOOK\_URL} and any local
overrides. If \texttt{.venv/} is missing, run
\texttt{python3 -m venv .venv} before the
\texttt{source} step.}

\newthought{Run the catch-up on \texttt{prod}.}
Same sequence on prod, where \texttt{setup-server.sh} also restarts the
service so the new logging, health, and analytics code goes live. After
the script finishes, verify the API is reachable from \texttt{dev}.

\marginnote{\newthought{Sequence on \texttt{prod}.}\\[0.4em]
\texttt{ssh produser@192.168.56.11}\\
\texttt{cd /opt/pixelwise}\\
\texttt{git fetch origin -{}-tags}\\
\texttt{git reset -{}-hard v0.8}\\
\texttt{cp .env.example .env}\\
\texttt{nano .env}\\
\texttt{source .venv/bin/activate}\\
\texttt{pip install -r requirements.txt}\\
\texttt{bash setup-server.sh}\\[0.3em]
Then from \texttt{dev}:
\texttt{curl http://192.168.56.11:8000/health}
should return \texttt{\{"status":"ok"\dots\}}.}

\marginnote{\newthought{Already have a fork?}\index{Git!fork}
If your own fork carries the \texttt{v0.8} tag, swap \texttt{origin} on
\texttt{dev} with
\texttt{git remote set-url origin git@github.com:<you>/pixelwise.git}
so push and pull continue to target your account.
The HTTPS clone on \texttt{prod} stays read-only.}

\newthought{You are now} at the state where Block~9 ended, with both
machines aligned, ready to start Block~10.

% ==============================================================================================
% BLOCK 2: CONCEPTS AND PRACTICE
% ==============================================================================================

\newpage
\section{The Agent Loop}\index{agent loop}

\newthought{How these tools actually work} the agent loop repeats until the task is done:

\begin{itemize}
    \item Observe: the agent collects context, your repo structure, git status, recent commits,
    the files you point it at.
    \item Inspect: it analyses the code and your request to understand what needs to change.
    \item Choose: it selects an action, edit a file, run a command, ask you a question.
    \item Act: it executes the action, gets the result, and loops back to observe.
\end{itemize}

\newthought{The harness} (the program wrapping the LLM\index{LLM}) validates every action before execution:
is the tool recognised? Are the arguments valid? Is the file path inside the repo? The LLM proposes,
the harness gatekeeps.

\newthought{The quality depends entirely on what context you give the agent.} Garbage in, garbage out.
A vague prompt produces vague code. A precise prompt with file references produces precise changes.

\subsection{Sandboxing}\index{sandboxing}

\newthought{AI agents run on your machine,} so what stops them from doing damage? An unrestricted agent
can run \texttt{rm -rf /}, \texttt{curl} your secrets to an external server, install packages, or modify
system files. It runs with your user's permissions.

\newthought{How sandboxing works in practice:}

\begin{description}
    \item[Permission prompts] Agents like Claude Code ask before executing each command, you approve or deny.
    The human stays in the loop.
    \item[Code-only mode] Aider is code-only by default, it reads and edits files but does not run shell
    commands at all.
    \item[Restricted execution] Docker containers, VMs, or chroot jails, the agent runs in an isolated
    environment where it physically cannot access the host system.
    \item[Allow/deny lists] Configure which commands or directories the agent can access. Some tools let
    you whitelist specific tools (e.g., ``can run \texttt{pytest} but nothing else'').
\end{description}

\marginnote{The principle: give the agent the minimum access it needs. Read your code, yes. Edit files, yes.
Run arbitrary commands, only with explicit approval. Access \texttt{.env} or SSH keys, never.}

% ==============================================================================================
\newpage
\section{Exercise: Set up Aider}\index{exercises}

\newthought{Set up Aider} on your dev VM. Aider\index{Aider} (\url{https://aider.chat/}) is an open-source CLI agent. Connect it to
Qwen\index{Qwen} (\url{https://qwen.readthedocs.io/}) via Ollama\index{Ollama} (\url{https://ollama.com/}), self-hosted on the server, provided by the course.
No external API needed:

\begin{verbatim}
pip install aider-chat
aider --model ollama/qwen2.5-coder
\end{verbatim}

\newthought{Once Aider runs}, you have a working instance of the agent loop sitting at your prompt. The next exercises will use the same session as their starting point.

% ==============================================================================================
\newpage
\section{Git as a Safety Net}\index{Git!safety net}

\newthought{Commit before you let an agent touch your code.} This is not just an AI workflow, it is how
you work with any risky change: experimental refactors, dependency upgrades, merge conflict resolution.

\begin{verbatim}
# Save current state
git add -A && git commit -m "pre-agent checkpoint"

# See what the agent changed
git diff

# Revert everything if the agent made a mess
git checkout .

# Stash your current state before experimenting
git stash

# Let the agent work on a branch
git checkout -b experiment
# ... agent does its work ...
# If it works, merge. If not, delete the branch.
git checkout main
git branch -D experiment
\end{verbatim}

\marginnote{Always commit clean state first, then change, then review the diff. This pattern predates AI
tools, it is how experienced developers have always worked with risky changes.}

\subsection{Reading Diffs}\index{diff}\index{code review}

\newthought{The core review skill.} When an agent (or a colleague) makes changes, you do not re-read
the entire file. You read the diff: what was added, what was removed, what was modified.

\begin{verbatim}
git diff                    # in the terminal
\end{verbatim}

\newthought{\texttt{git diff} in the terminal,} side-by-side diffs in the IDE, GitHub's pull request diff
view, all the same skill. Students need to read diffs fluently. This is how code review actually works.

\marginnote{Lines starting with \texttt{+} were added. Lines starting with \texttt{-} were removed.
Lines starting with a space are unchanged context. The \texttt{@@} header tells you which file and which
line numbers changed.}

% ==============================================================================================
\newpage
\section{Exercise: Commit Clean State}\index{exercises}

\newthought{First rule, commit clean state} before letting the agent touch anything:

\begin{verbatim}
git add -A
git commit -m "pre-agent checkpoint"
\end{verbatim}

This is your undo button for everything that follows. From this clean state, every change the agent makes shows up cleanly in \texttt{git diff} and can be reverted with a single command.

\newthought{Verify the safety net} works. Let the agent make a small, deliberately throwaway change, then read the diff and revert:

\begin{verbatim}
git diff
git checkout .
\end{verbatim}

You should land back at the pre-agent checkpoint with the agent's edits gone. Until you can do this without thinking, do not give the agent anything important to touch.

% ==============================================================================================
\newpage
\section{Context Management}\index{context window}\index{context management}

\newthought{Agents have a limited context window}, how much text the LLM can process at once.
Long conversations degrade quality: the agent forgets earlier instructions, repeats mistakes, or
contradicts itself. Older parts of the conversation get compressed or dropped.

\newthought{Practical rules:}
\begin{itemize}
    \item Start a fresh conversation when switching tasks.
    \item Point the agent at specific files rather than letting it read everything.
    \item Keep instructions files short and focused.
\end{itemize}

\subsection{Subagents}\index{subagent}

\newthought{Some tools can spawn child agents for subtasks}, e.g., ``research this question in a
separate context, report back.'' The child agent operates with tighter permissions and its own context
window, keeping the parent conversation clean.

% ==============================================================================================
\newpage
\section{Exercise: Explore the Tool}\index{exercises}

\newthought{Get comfortable with the conversation loop} by asking the agent to explain a file, describe the PixelWise architecture, or find where a function is called. Try:

\begin{itemize}
    \item ``Explain what \texttt{app/classifier.py} does''
    \item ``Where is \texttt{classify\_batch} called?''
    \item ``Describe the full request flow from browser to database''
\end{itemize}

\newthought{Watch the context fill up.} Notice how the agent's answers degrade if you keep piling questions into the same conversation, and how a fresh session with a single, file-scoped prompt produces a sharper answer. Pointing the agent at specific files rather than asking it to ``look at the whole project'' is the cheapest quality gain you have.

% ==============================================================================================
\newpage
\section{Teaching the AI Your Project}\index{instructions file}

\newthought{Instructions as code}, declarative files that persist across sessions. You write them once,
the AI follows them every time. This is the highest-leverage thing you can do with AI tools.

\begin{description}
    \item[\texttt{CLAUDE.md}]\index{CLAUDE.md} Project-level instructions for Claude Code: conventions,
    architecture, what NOT to do.
    \item[\texttt{.github/copilot-instructions.md}] Same concept for GitHub Copilot.
    \item[\texttt{.cursorrules}]\index{.cursorrules} Cursor-specific project rules.
\end{description}

\newthought{The pattern across all tools:} each has its own filename, but the idea is identical, a
markdown file in your repo that the AI reads before starting work.

\subsection{Skill Files}\index{skill file}

\newthought{Skill files are reusable recipes} for specific tasks, structured \texttt{.md} files with
step-by-step instructions the AI follows consistently. Example: ``how to add a new API endpoint to
this project.'' The AI produces correct code when given only the skill file and a task description.

% ==============================================================================================
\newpage
\section{Exercise: Write a CLAUDE.md}\index{exercises}

\newthought{Write a \texttt{CLAUDE.md}} (or equivalent instructions file) for the PixelWise project.
Document: the stack, the conventions, what the AI should and should not do. Example structure:

\begin{verbatim}
# PixelWise

## Stack
- Python 3.11, FastAPI, SQLAlchemy, PostgreSQL
- Nginx reverse proxy, vanilla JS frontend
- sklearn LogisticRegression model (MNIST)

## Conventions
- All config via .env (never hardcode secrets)
- Batch-first inference (classify_batch is the core)
- Pydantic models for all API contracts

## Do NOT
- Modify .env or expose secrets
- Add new dependencies without updating requirements.txt
- Change the classifier interface without updating tests
\end{verbatim}

Commit it to the repo. From now on, every agent session in this project starts with this file already in context.

% ==============================================================================================
\newpage
\section{Exercise: Write a Skill File}\index{exercises}

\newthought{Create a \texttt{.skills/} directory} and write a skill for a
specific task, e.g., ``how to add a new API endpoint'':

\begin{verbatim}
mkdir -p .skills
nano .skills/add-endpoint.md
\end{verbatim}

Include: which files to modify, what patterns to follow, what to test. Give the skill file to
the agent with a task description, does it produce correct code?

\newthought{Iterate on the skill} until the agent can carry out the task from the skill alone, without you correcting it mid-way. A good skill file is one a tired colleague could also follow.

% ==============================================================================================
\newpage
\section{Working with AI Agents Effectively}

\newthought{The agent-edit-verify loop:}
\begin{itemize}
    \item Agent makes changes.
    \item \texttt{git diff}, review every line.
    \item Run the app, works or does not.
    \item If not, feed the error back to the agent or fix it yourself.
    \item Iterate until it is right, or revert and try a different approach.
\end{itemize}

\newthought{Use \texttt{/plan}} to let the agent think before it acts, review the plan, adjust, then
let it execute. Do not jump straight to implementation.

\newthought{Code review with AI:} have the agent review your own code, it catches things you are blind
to after staring at the same file for hours.

% ==============================================================================================
\newpage
\section{Exercise: Make a Real Change with the AI}\index{exercises}

\newthought{Student's choice}, a new feature, a refactor, a bug fix.
Follow the agent-edit-verify loop:

\begin{itemize}
    \item Let the agent make changes.
    \item Review: \texttt{git diff}, read every line.
    \item Run the app, does it work?
    \item If not, feed the error back. Iterate or revert.
\end{itemize}

\newthought{Stay in the loop} the entire time. The moment you stop reading the diffs is the moment the agent starts shipping bugs into your repo with your name on the commit.

% ==============================================================================================
\newpage
\section{Exercise: Branching Workflow}\index{exercises}

\newthought{Create a branch for the AI's work} so the main line stays clean while you experiment:

\begin{verbatim}
git checkout -b ai-experiment
# ... let the agent work ...
git diff main
# If good: merge. If bad: delete the branch.
\end{verbatim}

The branch is the boundary between ``the agent has free rein'' and ``the codebase you trust''. Merge only what passes the diff and the run.

% ==============================================================================================
\newpage
\section{Exercise: Group Discussion}\index{exercises}

\newthought{Compare notes with the room.} When was the AI faster than writing code by hand? When was it slower?
When was it confidently wrong? What did the instructions file change about the AI's output quality?

\newthought{Name one thing} you will keep doing, one thing you will stop doing, and one habit you want to acquire before the agent touches production code on your behalf.

% ==============================================================================================
\newpage
\section{Optional: MCP and Tool Use}\index{MCP}\index{Model Context Protocol}

\newthought{Agents can do more than edit code.} MCP (Model Context Protocol)\marginnote{MCP\index{MCP}:
\url{https://modelcontextprotocol.io/}} lets agents connect to
external tools, reading documentation, searching the web, running tests, querying databases, interacting
with APIs. The agent is not limited to what is in your repo.

% ==============================================================================================
\newpage
\section{Optional: When to Trust AI Output}\index{AI trust}

\begin{description}
    \item[Strong] Boilerplate, tests for known behaviour, documentation, refactoring, explaining unfamiliar code.
    \item[Weak] Security-sensitive code, novel algorithms, anything touching infrastructure or secrets.
    \item[Never] Blind trust without reading the output.
\end{description}

% ==============================================================================================
\newpage
\section{Security Thread}

\newthought{AI-generated code must pass the same security review as human-written code.}
Prompt injection\index{prompt injection}: user-submitted text that ends up in an AI prompt is an injection
vector, never pipe untrusted input into an AI call. The AI must never be given access to \texttt{.env},
SSH keys, or production credentials directly.

\newthought{A note on what comes next.} This is the final block. The course wrap-up sits in the recap below: every layer you touched, from two empty VMs to a deployed, monitored, AI-assisted service, is now yours to maintain.

% ==============================================================================================
% BLOCK 3: CONSOLIDATION
% ==============================================================================================

\section{Self-Reflection and Recap}

\newthought{Reflection questions:}
\begin{itemize}
    \item Why do you commit clean state before letting an agent touch your code?
    \item What is the difference between a CLI agent and an IDE agent? When would you use each?
    \item How does an instructions file improve the quality of AI-generated code?
    \item What are the risks of giving an AI agent unrestricted access to your machine?
    \item When should you trust AI output, and when should you be sceptical?
\end{itemize}

\newthought{Recap} of key concepts:
\begin{itemize}
    \item The agent loop (observe, inspect, choose, act) and the harness that gatekeeps every action.
    \item Git as the safety net: clean checkpoint, branch, diff, revert.
    \item Sandboxing and least privilege: the agent gets the minimum access it needs.
    \item Context management: short conversations, file-scoped prompts, instructions files.
    \item \texttt{CLAUDE.md} and skill files as the highest-leverage way to steer AI output.
\end{itemize}

\marginnote{\newthought{Milestone.} Students have used a terminal AI agent on a real codebase. Instructions file and
skill file committed to the repo. Every student can read a diff, revert a bad change, and iterate with an agent.}

\marginnote{\newthought{Course wrap-up.} Look at what you have built, from two empty VMs to a deployed, monitored,
auto-deploying AI web service with structured logging, database persistence, CI/CD, and now AI-assisted
development. The full stack, end to end. You own every layer.}
